\documentclass[11pt,a4paper]{article}
\usepackage[T1]{fontenc}
\usepackage[utf8]{inputenc}
\usepackage{lmodern}
\usepackage[a4paper,margin=25mm,headheight=14pt]{geometry}
\usepackage{amsmath,amssymb,amsthm,mathtools}
\usepackage{booktabs,array,enumitem,microtype,needspace}
\usepackage[dvipsnames]{xcolor}
\usepackage{fancyhdr}
\usepackage{hyperref}
\hypersetup{colorlinks=true,linkcolor=MidnightBlue,citecolor=MidnightBlue,urlcolor=MidnightBlue,pdftitle={Gamma and q-Gamma Kernels: Global Determinant Signs and Applications},pdfauthor={Keio and Shaitetsu}}
\pagestyle{fancy}
\fancyhf{}
\fancyhead[L]{\small Gamma and q-gamma kernels}
\fancyhead[R]{\small Global signs and applications}
\fancyfoot[C]{\thepage}
\setlength{\parskip}{0.35em}
\setlength{\parindent}{1.2em}
\setlist{itemsep=0.25em,topsep=0.4em}
\allowdisplaybreaks[2]
\newtheorem{theorem}{Theorem}[section]
\newtheorem{lemma}[theorem]{Lemma}
\newtheorem{proposition}[theorem]{Proposition}
\newtheorem{corollary}[theorem]{Corollary}
\theoremstyle{definition}
\newtheorem{definition}[theorem]{Definition}
\newtheorem{example}[theorem]{Example}
\theoremstyle{remark}
\newtheorem{remark}[theorem]{Remark}
\DeclareMathOperator{\sgn}{sgn}
\DeclareMathOperator{\rank}{rank}
\DeclareMathOperator{\diag}{diag}
\newcommand{\R}{\mathbb R}
\newcommand{\N}{\mathbb N}
\newcommand{\Z}{\mathbb Z}
\newcommand{\dd}{\,\mathrm d}
\newcommand{\ep}{\varepsilon}
\newcommand{\V}{\mathcal V}
\newcommand{\SR}{\mathrm{SR}}
\newcommand{\SSR}{\mathrm{SSR}}
\newcommand{\STP}{\mathrm{STP}}
\newcommand{\simplex}[1]{\mathcal S_{#1}}

\title{\textbf{Gamma and \(q\)-Gamma Kernels:}\\[0.25em]
\textbf{Global Determinant Signs and Applications}\\[0.5em]
\large Reciprocal kernels, complete gamma-ratio classification,\\
product obstructions, and shape constraints}
\author{Keio and Shaitetsu}
\date{20 November 2026}

\begin{document}
\maketitle
\begin{abstract}
We present a unified treatment of determinant signs for reciprocal
Jackson \(q\)-gamma, reciprocal gamma, and gamma-ratio kernels on
their positive-argument domains. Both reciprocal kernels have strict
signature \((-1)^{r(r-1)/2}\) at every finite order; the ordinary
reciprocal gamma function also satisfies the corresponding strict
derivative Hankel inequalities. For
\(\Gamma(c+x+y)/\Gamma(d+x+y)\), the signature is determined
completely by \(\prod_{j=0}^{r-1}(d-c)_j\). A nonnegative integer
difference \(c-d=m\) gives exact rank \(m+1\), while every other
real difference gives strict signs at all orders. The proofs use
discrete and continuous determinant composition, a direct beta
representation, and a gamma-density transform of a power kernel.
We also exhibit an exact family of strictly sign-regular factors whose
pointwise product has second-order minors of both signs. Applications
include generalized Descartes bounds, unique interpolation, monotone
ratios, an explicit unimodality transition, and convergent infinite
series. All determinant and shape arguments are proved within the
paper. For a separate nonlinear question about reciprocal \(q\)-gamma
sums, we prove a curvature identity and identify the additional
variance bound that is still required; that bound is not claimed as
a theorem.
\end{abstract}

\noindent\textbf{Keywords.} Sign-regular kernel; total positivity;
\(q\)-gamma function; gamma ratio; Cauchy--Binet identity;
variation of signs; unimodality; Tur\'an inequality.

\section{Introduction and scope}

This paper studies three positive kernels whose determinants have a
common global structure:
\[
 \frac{1}{\Gamma_q(x+y)},\qquad
 \frac{1}{\Gamma(x+y)},\qquad
 K_{c,d}(x,y)=\frac{\Gamma(c+x+y)}{\Gamma(d+x+y)}.
\]
The first two have one universal alternating signature. The third has
a signature depending on \(c-d\), with exact finite-rank degeneracies.
We give their proofs in a single framework and develop consequences for
zeros, interpolation, ratios, and the preservation or reversal of
unimodality. The reciprocal-kernel questions and the gamma-ratio
assertion appear in \cite[Section~3, Examples~6, 8, and~11]{KVZ}.

Throughout, all nodes are strictly increasing unless otherwise stated.
The adjective \emph{global} means that the stated sign holds for every
finite order and every admissible pair of node lists; it does not refer
to a conclusion drawn from numerical tests. The natural domains require
positive gamma arguments. They permit real nodes of either sign when
the sums satisfy that condition.

The results and their logical status are summarized below. The last row
is an explicitly identified further question; it is not used in any
proof of the preceding results.

\begin{center}
\small
\begin{tabular}{@{}p{0.31\textwidth}p{0.59\textwidth}@{}}
\toprule
Object & Result established in this paper\\
\midrule
Reciprocal \(q\)-gamma & Every ordered minor has strict sign
\((-1)^{\binom r2}\); finite and inverse finite \(q\)-factorial
consequences follow (Theorem~\ref{qg:main}).\\[0.45em]
Reciprocal gamma & The same strict signs hold, including a strict
derivative Hankel hierarchy and real-order inverse factorials
(Theorems~\ref{cg:main} and~\ref{cg:confluent}).\\[0.45em]
Gamma ratios & The signature is
\(\sgn\prod_{j=0}^{r-1}(d-c)_j\); nonnegative integer
\(c-d\) gives exact finite rank (Theorem~\ref{thm:main}).\\[0.45em]
Pointwise products & Two strictly sign-regular factors can have a
product with second-order minors of both signs
(Theorem~\ref{pc:counterexample}).\\[0.45em]
Linear and ratio applications & Distinct-zero bounds, exact
interpolation, monotonicity, and unimodality follow under the stated
rank and convergence conditions (Section~\ref{app:section}).\\[0.45em]
Nonlinear \(q\)-gamma sums & An exact curvature identity is proved.
The further variance bound needed for continuous-parameter
log-concavity remains unproved here (Section~\ref{ob:nonlinear}).\\
\bottomrule
\end{tabular}
\end{center}

The determinant proofs are supplied in full, including the strictness
arguments in both series and integral limits. The unimodality
consequence is also explained below. This manuscript makes no claim to
have exhausted the publication history of each statement.

\section{Determinant tools and notation}

Throughout, \((a)_0=1\) and
\((a)_j=a(a+1)\cdots(a+j-1)\) denote rising factorials. Put
\begin{equation}\label{eq:sign-product}
 C_r(\delta)=\prod_{j=0}^{r-1}(-\delta)_j,
 \qquad \ep_r=(-1)^{\binom r2}.
\end{equation}
The product definition, rather than a gamma quotient, is used when a factor
vanishes. For an increasing tuple \(z_1<\cdots<z_r\), write
\(\V(z)=\prod_{i<j}(z_j-z_i)>0\).

\begin{definition}
A kernel on a product of ordered sets is strictly sign regular through
order \(R\) if, for each \(1\le r\le R\), all its ordered \(r\)-minors
are nonzero with one common sign \(\sigma_r\in\{+1,-1\}\).
Strict total positivity means \(\sigma_r=+1\). Non-strict sign regularity
permits zero minors. The subscript \(\infty\) means every finite order.
\end{definition}

\subsection{From Wronskians to ordered determinants}


We use the Wronskian convention
\[
 W(f_1,\ldots,f_k)(u)
 =\det[f_j^{(i-1)}(u)]_{i,j=1}^{k}.
\]

\begin{lemma}[A Rolle--Wronskian criterion]\label{lem:rolle}
Let \(f_1,\ldots,f_r\) be smooth real functions on an open interval \(I\).
Suppose every prefix Wronskian
\(W(f_1,\ldots,f_k)\), \(1\le k\le r\), is nowhere zero on \(I\).
Then every nonzero linear combination has at most \(r-1\) distinct zeros
on \(I\). For every \(u_1<\cdots<u_r\) in \(I\),
\(\det[f_j(u_i)]\) is nonzero and has the sign of
\(W(f_1,\ldots,f_r)\).
\end{lemma}

\begin{proof}
Induct on \(r\). The case \(r=1\) is immediate. Since \(f_1\) never
vanishes, define \(q_j=(f_j/f_1)'\), \(2\le j\le r\). The Wronskian
product identity and expansion along the first column give
\begin{equation}\label{eq:quotient-wronskian}
 W(q_2,\ldots,q_k)
 =\frac{W(f_1,\ldots,f_k)}{f_1^k},\qquad 2\le k\le r.
\end{equation}
To verify the identity, first use
\(W(f_1g_1,\ldots,f_1g_k)=f_1^k W(g_1,\ldots,g_k)\), which follows
from the triangular Leibniz rule, and then use
\(W(1,g_2,\ldots,g_k)=W(g_2',\ldots,g_k')\).

If a nonzero combination \(F=\sum_j a_jf_j\) had \(r\) distinct zeros,
Rolle's theorem would give at least \(r-1\) distinct zeros of
\((F/f_1)'=\sum_{j=2}^r a_jq_j\). If some \(a_j\), \(j\ge2\), is
nonzero, the induction hypothesis applied using
\eqref{eq:quotient-wronskian} gives the contradictory upper bound
\(r-2\). Such a combination cannot vanish identically, since its
Wronskian is nonzero. If instead all those coefficients vanish, then
\(F=a_1f_1\) has no zeros. This proves the zero bound and consequently
the nonvanishing of every collocation determinant.

The increasing-node domain is connected, so the continuous nonvanishing
determinant has one sign on that domain. For coalescing increasing nodes,
Taylor expansion gives
\begin{equation}\label{eq:coalescent}
 \lim_{u_1,\ldots,u_r\to u}
 \frac{\det[f_j(u_i)]}{\V(u_1,\ldots,u_r)}
 =\frac{W(f_1,\ldots,f_r)(u)}{\prod_{j=0}^{r-1}j!}.
\end{equation}
One may take \(u_i=u+(i-1)h\) with \(h\downarrow0\). The denominator
on the right and the Vandermonde on the left are positive. This determines
the sign.
\end{proof}

\subsection{Exponential and generalized Vandermonde kernels}

\begin{lemma}[Generalized Vandermonde signs]\label{lem:monomials}
For arbitrary real \(a_1<\cdots<a_r\) and positive
\(t_1<\cdots<t_r\),
\begin{equation}\label{cg:generalized-vandermonde}
 \det[t_i^{a_j}]_{i,j=1}^r>0.
\end{equation}
If the positive nodes are strictly decreasing, the strict sign is
\(\ep_r\). Transposition does not alter either conclusion.
In particular, the exponential kernel \((u,v)\mapsto e^{uv}\)
is strictly totally positive on \(\R\times\R\).
\end{lemma}
\begin{proof}
For increasing real exponents, every prefix Wronskian satisfies
\[
 W(e^{a_1z},\ldots,e^{a_kz})
 =e^{z\sum_{j=1}^k a_j}\V(a_1,\ldots,a_k)>0.
\]
Lemma~\ref{lem:rolle} proves positivity of the exponential
collocation determinants. Set \(z_i=\log t_i\) to obtain the
power determinant. Reversing \(r\) rows multiplies a determinant
by \((-1)^{\binom r2}\), proving the decreasing-node statement.
\end{proof}

\subsection{The power-kernel signature}

\begin{proposition}[The power kernel]\label{prop:power}
For \(u,v>0\), let \(P_\delta(u,v)=(u+v)^\delta\).
If \(\delta\notin\Z_{\ge0}\), this kernel is strictly sign regular of
every order, with signature \(\sgn C_r(\delta)\).
If \(\delta=m\in\Z_{\ge0}\), it has strict sign \(\ep_r\) through
order \(m+1\), and all higher minors vanish.
\end{proposition}

\begin{proof}
Fix \(0<v_1<\cdots<v_r\) and put \(g_j(u)=(u+v_j)^\delta\).
For each \(1\le k\le r\), direct differentiation gives
\begin{equation}\label{eq:power-wronskian}
 W(g_1,\ldots,g_k)(u)
 =C_k(\delta)\,\V(v_1,\ldots,v_k)
   \prod_{j=1}^{k}(u+v_j)^{\delta-k+1}.
\end{equation}
Here is the sign calculation explicitly. Factor
\(\delta^{\underline i}=\delta(\delta-1)\cdots(\delta-i+1)\)
from derivative row \(i\), and factor \((u+v_j)^\delta\) from column
\(j\). The remaining determinant is the Vandermonde on the decreasing
numbers \(1/(u+v_j)\), hence equals
\[
 (-1)^{\binom k2}\V(v_1,\ldots,v_k)
 \prod_{j=1}^{k}(u+v_j)^{-(k-1)}.
\]
On the other hand,
\(\prod_{i=0}^{k-1}\delta^{\underline i}
 =(-1)^{\binom k2}C_k(\delta)\).
The two signs cancel, proving \eqref{eq:power-wronskian}.

If \(\delta\notin\Z_{\ge0}\), all these prefix Wronskians are
nonzero, and Lemma~\ref{lem:rolle} proves the assertion. If
\(\delta=m\ge0\) is an integer, the same argument applies for
\(k\le m+1\), where \(\sgn C_k(m)=\ep_k\). For higher orders,
\[
 (u+v)^m=\sum_{j=0}^{m}\binom mj u^jv^{m-j}
\]
has rank at most \(m+1\), so the corresponding minors vanish.
\end{proof}

\subsection{Finite and continuous determinant composition}

\begin{lemma}[Finite Cauchy--Binet]\label{cb:finite}
Let \(A\) be an \(r\times M\) matrix and \(B\) an
\(M\times r\) matrix, with \(M\ge r\). Then
\begin{equation}\label{cb:finite-equation}
 \det(AB)=
 \sum_{1\le j_1<\cdots<j_r\le M}
 \det[A_{i,j_k}]_{i,k=1}^r
 \det[B_{j_k,l}]_{k,l=1}^r.
\end{equation}
\end{lemma}
\begin{proof}
Expand the determinant of \(AB\) by multilinearity in its columns.
Any term choosing the same column of \(A\) twice vanishes.
For a fixed set of \(r\) distinct chosen indices, the sum over their
orders is the determinant of the corresponding rows of \(B\),
multiplied by the determinant of the corresponding columns of \(A\).
Grouping the terms by their increasing index sets gives the formula.
\end{proof}

\begin{lemma}[Ordered integral Cauchy--Binet]\label{cb:integral}
Let \(J\) be a real interval with a positive Borel measure \(\nu\).
Suppose measurable real functions \(A_i,B_j\), \(1\le i,j\le r\),
satisfy \(\int_J|A_i(t)B_j(t)|\,\mathrm d\nu(t)<\infty\) for every
pair. Put \(M_{ij}=\int_J A_i(t)B_j(t)\,\mathrm d\nu(t)\).
Then
\begin{align}\label{cb:integral-equation}
 \det M
 &=\frac1{r!}\int_{J^r}
  \det[A_i(t_k)]_{i,k=1}^r
  \det[B_j(t_k)]_{k,j=1}^r
  \prod_{k=1}^r\mathrm d\nu(t_k)\notag\\
 &=\int_{\substack{t_1<\cdots<t_r\\t_k\in J}}
  \det[A_i(t_k)]_{i,k=1}^r
  \det[B_j(t_k)]_{k,j=1}^r
  \prod_{k=1}^r\mathrm d\nu(t_k).
\end{align}
\end{lemma}
\begin{proof}
Expanding the two determinants gives finitely many terms. Each absolute
term is a product of \(r\) integrable factors of the form
\(|A_i(t_k)B_j(t_k)|\), so Fubini applies to the full expansion.
After relabeling the variables, every permutation in the first
determinant gives the same signed expansion of \(\det M\).
There are \(r!\) such permutations, proving the first equality.
The integrand is invariant under simultaneous permutation of the
variables, because both determinants change by the same sign.
If two variables coincide, both determinants vanish. Thus the
\(r!\) strict ordering regions exhaust its contribution, proving the
second equality, even when the measure has atoms.
\end{proof}

In a finite truncation of a series composition, strictness survives the
limit if one fixed summand in \eqref{cb:finite-equation} has a strictly
positive signed contribution independent of the truncation. In an
integral composition, a continuous integrand with a strict sign on a
nonempty open set of positive measure likewise gives a strict
determinant. These are the strictness mechanisms used below.

\section{Reciprocal Jackson \texorpdfstring{$q$}{q}-gamma and finite
\texorpdfstring{$q$}{q}-factorials}\label{qg:section}

Fix \(0<q<1\). For \(n\in\mathbb Z_{\ge0}\), write
\[
 (u;q)_n=\prod_{j=0}^{n-1}(1-uq^j),\qquad
 (u;q)_\infty=\prod_{j=0}^{\infty}(1-uq^j),
 \qquad (u;q)_0=1.
\]
For \(0\le u<1\), the infinite product is strictly positive: its
factors are positive and the sum of their absolute logarithms is finite.
For \(s>0\), Jackson's gamma function is
\begin{equation}\label{qg:definition}
 \Gamma_q(s)=(1-q)^{1-s}
 \frac{(q;q)_\infty}{(q^s;q)_\infty}>0.
\end{equation}
Splitting off the first factor of the infinite product gives
\begin{equation}\label{qg:recurrence}
 \Gamma_q(1)=1,\qquad
 \Gamma_q(s+1)=\frac{1-q^s}{1-q}\,\Gamma_q(s),\qquad s>0.
\end{equation}
As elsewhere, put \(\ep_r=(-1)^{\binom r2}\).

\begin{theorem}[Reciprocal Jackson gamma kernel]\label{qg:main}
For every \(0<q<1\), every \(r\ge1\), and arbitrary increasing
real node lists \(x_1<\cdots<x_r\), \(y_1<\cdots<y_r\)
satisfying \(x_1+y_1>0\),
\begin{equation}\label{qg:main-sign}
 \ep_r\det\left[\frac{1}{\Gamma_q(x_i+y_j)}\right]_{i,j=1}^r>0.
\end{equation}
Thus the reciprocal kernel is strictly sign regular of every order on
every product domain on which all the sums are positive.
\end{theorem}

The proof supplies an explicit factorization and a positive lower
bound for each signed determinant before passing to an infinite-series
limit.

\subsection{The binomial identity}

\begin{lemma}[The \(q\)-binomial identity]\label{qg:binomial}
For \(0\le A<1\) and \(|z|<1\),
\begin{equation}\label{qg:binomial-equation}
 \frac{(Az;q)_\infty}{(z;q)_\infty}
 =\sum_{n=0}^{\infty}\frac{(A;q)_n}{(q;q)_n}z^n.
\end{equation}
The products and the series converge, and the series converges
absolutely and uniformly on every closed disk \(|z|\le\rho<1\).
\end{lemma}
\begin{proof}
For \(|w|\le\rho<1\), the power-series branch of the logarithm
satisfies
\[
 |\log(1-w)|\le \frac{|w|}{1-\rho}.
\]
Consequently the logarithmic series defining \((z;q)_\infty\)
and \((Az;q)_\infty\) converge uniformly on every such disk,
because \(\sum_{j\ge0}q^j<\infty\). Their exponentials are analytic
and nonvanishing there. Thus
\[
 F(z)=\frac{(Az;q)_\infty}{(z;q)_\infty}
\]
is analytic for \(|z|<1\), with \(F(0)=1\). Splitting off the
initial product factors yields
\[
 (1-z)F(z)=(1-Az)F(qz).
\]
If \(F(z)=\sum_{n\ge0}c_nz^n\), comparison of coefficients gives
\[
 c_0=1,\qquad
 (1-q^n)c_n=(1-Aq^{n-1})c_{n-1}\quad(n\ge1).
\]
Hence \(c_n=(A;q)_n/(q;q)_n\), proving the identity.
For the stated convergence directly, note that
\[
 0<\frac{(A;q)_n}{(q;q)_n}
 \le\frac{1}{(q;q)_\infty},
\]
so a geometric series dominates the right side on each closed disk.
\end{proof}

\subsection{Total positivity of the finite shifted-factorial kernel}

\begin{theorem}[Finite \(q\)-factorials at every order]\label{qg:factorials}
For every \(r\ge1\), \(0<x_1<\cdots<x_r\), and
\(0\le n_1<\cdots<n_r\) with integer \(n_j\),
\begin{equation}\label{qg:factorial-sign}
 \det[(q^{x_i};q)_{n_j}]_{i,j=1}^r>0.
\end{equation}
Thus \((x,n)\mapsto(q^x;q)_n\) is \(\STP_\infty\) on
\((0,\infty)\times\mathbb Z_{\ge0}\).
\end{theorem}
\begin{proof}
The case \(A=0\) of Lemma~\ref{qg:binomial} gives
\begin{equation}\label{qg:auxiliary-expansion}
 H(x,n):=\frac{1}{(q^{x+n};q)_\infty}
 =\sum_{k=0}^{\infty}
   \frac{(q^x)^k(q^n)^k}{(q;q)_k}.
\end{equation}
This converges since \(q^{x+n}<1\). Put \(u_i=q^{x_i}\) and
\(v_j=q^{n_j}\). Both node lists are strictly decreasing and
positive; the possible endpoint \(v_1=1\) causes no difficulty.
Let \(H^{(M)}\) denote the matrix obtained by truncating
\eqref{qg:auxiliary-expansion} at \(k=M\).
For \(M\ge r-1\), finite Cauchy--Binet
(Lemma~\ref{cb:finite}) gives
\begin{equation}\label{qg:first-cb}
 \det H^{(M)}
 =\sum_{0\le k_1<\cdots<k_r\le M}
   \frac{
    \det[u_i^{k_h}]_{i,h=1}^r
    \det[v_j^{k_h}]_{h,j=1}^r}
   {\prod_{h=1}^r(q;q)_{k_h}}.
\end{equation}
Each of the two determinants in every numerator has strict sign
\(\ep_r\), by Lemma~\ref{lem:monomials}; their product is
strictly positive. In particular, the single index set
\(\{k_1,\ldots,k_r\}=\{0,\ldots,r-1\}\) contributes a fixed
number \(L_H>0\), independent of \(M\). Therefore
\[
 \det H^{(M)}\ge L_H>0\qquad(M\ge r-1).
\]
The entries converge to those of \(H\), so continuity of the
finite-dimensional determinant gives \(\det H\ge L_H>0\).
Finally,
\[
 (q^x;q)_n=(q^x;q)_\infty H(x,n).
\]
The strictly positive row factors \((q^{x_i};q)_\infty\)
transfer this strict positivity to \eqref{qg:factorial-sign}.
\end{proof}

\subsection{The binomial bridge to the reciprocal kernel}

\begin{proof}[Proof of Theorem~\ref{qg:main}]
First suppose both node lists are strictly positive. Set
\(z_j=q^{y_j}\), so \(1>z_1>\cdots>z_r>0\). Apply
Lemma~\ref{qg:binomial} with \(A=q^{x_i}\) to obtain
\begin{equation}\label{qg:bridge}
 T_{ij}:=\frac{(q^{x_i}z_j;q)_\infty}{(z_j;q)_\infty}
 =\sum_{n=0}^{\infty}
   (q^{x_i};q)_n\frac{z_j^n}{(q;q)_n}.
\end{equation}
Let \(T^{(M)}\) use only the indices \(0\le n\le M\).
For \(M\ge r-1\), Lemma~\ref{cb:finite} yields
\begin{align}\label{qg:second-cb}
 \ep_r\det T^{(M)}
 &=\sum_{0\le n_1<\cdots<n_r\le M}
   \det[(q^{x_i};q)_{n_h}]_{i,h=1}^r
   \frac{\ep_r\det[z_j^{n_h}]_{h,j=1}^r}
        {\prod_{h=1}^r(q;q)_{n_h}}.
\end{align}
The first determinant in each summand is strictly positive by
Theorem~\ref{qg:factorials}. The signed second determinant is
strictly positive by Lemma~\ref{lem:monomials}. Thus every summand
is positive, and the fixed index set \(\{0,\ldots,r-1\}\)
contributes a number \(L_T>0\) independent of \(M\). We have
\[
 \ep_r\det T^{(M)}\ge L_T>0\qquad(M\ge r-1).
\]
The series in \eqref{qg:bridge} converges absolutely, so its entries
converge to those of \(T\). It follows that
\(\ep_r\det T\ge L_T>0\). In particular, strictness survives
the infinite limit because of this fixed positive contribution;
it is not inferred merely from strictness of the truncations.

By \eqref{qg:definition}, the target entries are
\begin{equation}\label{qg:row-column-scaling}
 \frac{1}{\Gamma_q(x_i+y_j)}
 =(1-q)^{x_i}\,T_{ij}\,
   \frac{(1-q)^{y_j-1}(z_j;q)_\infty}{(q;q)_\infty}.
\end{equation}
These row and column multipliers are all strictly positive. Hence
\eqref{qg:main-sign} holds when \(x_i,y_j>0\).

For the full stated domain, choose a real number \(h\) with
\(-x_1<h<y_1\), which is possible precisely because
\(x_1+y_1>0\). The shifted lists \(x_i+h\) and \(y_j-h\)
are strictly positive and increasing, while every pairwise sum is
unchanged. The positive-node conclusion therefore proves
\eqref{qg:main-sign} for the original real nodes as well.
\end{proof}

\begin{corollary}[Finite inverse \(q\)-factorials at every order]
\label{qg:inverse-factorials}
For every \(r\ge1\), \(0<x_1<\cdots<x_r\), and increasing
nonnegative integers \(n_1<\cdots<n_r\),
\begin{equation}\label{qg:inverse-factorial-sign}
 \ep_r\det\left[\frac{1}{(q^{x_i};q)_{n_j}}\right]_{i,j=1}^r>0.
\end{equation}
Thus the reciprocal finite shifted-factorial kernel is
\(\SSR_\infty\), with signature \(\ep_r\).
\end{corollary}
\begin{proof}
Iterating \eqref{qg:recurrence} gives
\[
 \frac{1}{(q^x;q)_n}
 =(1-q)^{-n}\frac{\Gamma_q(x)}{\Gamma_q(x+n)}.
\]
Apply Theorem~\ref{qg:main} to the node lists \(x_i\) and
\(n_j\), then multiply row \(i\) by \(\Gamma_q(x_i)>0\)
and column \(j\) by \((1-q)^{-n_j}>0\). The extended sum domain
in that theorem explicitly permits \(n_1=0\), since
\(x_1+n_1>0\).
\end{proof}

\section{Ordinary reciprocal gamma and strict confluent inequalities}
\label{cg:section}

A direct beta-integral representation proves the reciprocal-gamma
assertion and strictness of every derivative Hankel determinant. These strict conclusions require more than entrywise
convergence in a limit of strictly signed determinants.

\begin{theorem}[Ordinary reciprocal-gamma kernel]\label{cg:main}
For every \(r\ge1\) and strictly increasing real node lists
\(x_1<\cdots<x_r\), \(y_1<\cdots<y_r\) satisfying \(x_1+y_1>0\),
\begin{equation}\label{cg:main-inequality}
 \ep_r\det\left[\frac{1}{\Gamma(x_i+y_j)}\right]_{i,j=1}^{r}>0,
 \qquad \ep_r=(-1)^{\binom r2}.
\end{equation}
In particular, \(K(x,y)=1/\Gamma(x+y)\) is \(\SSR_\infty\) on
\((0,\infty)\times(0,\infty)\), with signature \((\ep_r)_{r\ge1}\).
\end{theorem}

\begin{proof}
We first suppose that both node lists are positive. Euler's beta integral
is
\begin{equation}\label{cg:beta}
 B(x,y)=\int_0^1 t^{x-1}(1-t)^{y-1}\dd t
       =\frac{\Gamma(x)\Gamma(y)}{\Gamma(x+y)},\qquad x,y>0.
\end{equation}
For completeness, the identity follows by applying the substitution
\(s=u+v\), \(t=u/(u+v)\), whose Jacobian is \(s\), to the product of the
two positive gamma integrals. The resulting integral factors as
\(\Gamma(x+y)B(x,y)\).

The generalized Vandermonde inequality
\eqref{cg:generalized-vandermonde} permits arbitrary real exponents.

Write
\[
 \Omega_r=\{(t_1,\ldots,t_r):0<t_1<\cdots<t_r<1\}.
\]
The ordered continuous Cauchy--Binet formula
(Lemma~\ref{cb:integral}) applied to \eqref{cg:beta} gives
\begin{equation}\label{cg:beta-composition}
 \det[B(x_i,y_j)]_{i,j=1}^{r}
 =
 \int_{\Omega_r}
   \det[t_k^{x_i-1}]_{i,k=1}^{r}\,
   \det[(1-t_k)^{y_j-1}]_{k,j=1}^{r}
   \dd\mathbf t .
\end{equation}
Here and below \(\dd\mathbf t=\dd t_1\cdots\dd t_r\).
Absolute integrability can be checked before invoking the formula:
expanding the two determinants gives finitely many products of factors
\(t_k^{x_i-1}(1-t_k)^{y_j-1}\).
Every resulting full-cube integral is a product of finite beta
integrals. Thus the determinant expansion and Fubini's theorem are
legitimate. The ordered-domain formula has no additional factorial:
the two determinant factors change sign together under a permutation of
the \(t_k\), so the full-cube factor \(1/r!\) is exactly removed by
restriction to \(\Omega_r\).

By \eqref{cg:generalized-vandermonde}, the first determinant in
\eqref{cg:beta-composition} is strictly positive. The bases
\(1-t_1,\ldots,1-t_r\) are strictly decreasing, and reversing their
order gives
\[
 \ep_r\det[(1-t_k)^{y_j-1}]_{k,j=1}^{r}>0.
\]
Consequently, after multiplication by \(\ep_r\), the integrand in
\eqref{cg:beta-composition} is continuous and strictly positive at every
point of the nonempty open set \(\Omega_r\). Its integral is strictly
positive. Finally,
\begin{equation}\label{cg:beta-scaling}
 \det[B(x_i,y_j)]
 =
 \left(\prod_{i=1}^{r}\Gamma(x_i)\right)
 \left(\prod_{j=1}^{r}\Gamma(y_j)\right)
 \det\left[\frac1{\Gamma(x_i+y_j)}\right],
\end{equation}
and both scaling products are positive.

For the stated real-node domain, choose \(h\) with
\(-x_1<h<y_1\). Replacing \(x_i\) by \(x_i+h\) and \(y_j\) by \(y_j-h\)
makes both lists positive without changing any sum \(x_i+y_j\).
The already proved case establishes \eqref{cg:main-inequality}.
\end{proof}

\subsection{Strict derivative Hankel determinants}

Coalescence of the preceding collocation determinants would immediately
give weak derivative inequalities. The next theorem proves the strict
inequalities directly through a second positive integral.

\begin{theorem}[Strict confluent reciprocal-gamma inequalities]
\label{cg:confluent}
Let \(f(s)=1/\Gamma(s)\) for \(s>0\). For every \(r\ge1\) and \(s>0\),
\begin{equation}\label{cg:confluent-inequality}
 \ep_r H_r(s)>0,
 \qquad
 H_r(s)=\det[f^{(i+j)}(s)]_{i,j=0}^{r-1}.
\end{equation}
More precisely, choose any \(x,y>0\) with \(x+y=s\), and put
\(w_{x,y}(t)=t^{x-1}(1-t)^{y-1}\). Then
\begin{equation}\label{cg:confluent-integral}
\begin{split}
 H_r(s)
 ={}&\frac{1}{[\Gamma(x)\Gamma(y)]^r}
 \int_{\Omega_r}
   \left(\prod_{k=1}^{r}w_{x,y}(t_k)\right)
   \V(\log t_1,\ldots,\log t_r)\\
 &\hspace{32mm}\cdot
   \V(\log(1-t_1),\ldots,\log(1-t_r))
   \dd\mathbf t ,
\end{split}
\end{equation}
where for any real tuple \(z\) we write
\(\V(z_1,\ldots,z_r)=\prod_{i<j}(z_j-z_i)\).
\end{theorem}

\begin{proof}
For nonnegative integers \(i,j\), differentiation of
\eqref{cg:beta} under the integral gives
\begin{equation}\label{cg:beta-derivatives}
 M_{ij}(x,y):=\partial_x^i\partial_y^j B(x,y)
 =
 \int_0^1
   w_{x,y}(t)(\log t)^i(\log(1-t))^j\dd t .
\end{equation}
To justify all derivatives used here, fix a compact rectangle of positive
\((x,y)\)-values, and let its positive lower bounds be \(a,b\).
For each fixed pair \(i,j\), the absolute values of the differentiated
integrands are bounded by
\[
 t^{a-1}(1-t)^{b-1}
 \bigl(1+|\log t|+|\log(1-t)|\bigr)^{i+j}.
\]
This is integrable: near either endpoint, a positive power exponent
greater than \(-1\) absorbs every fixed logarithmic power. Dominated
convergence therefore justifies \eqref{cg:beta-derivatives}.

Apply the ordered composition formula to the matrix
\(M=(M_{ij})_{i,j=0}^{r-1}\), using the positive measure
\(w_{x,y}(t)\dd t\). The first determinant is formed from the powers
\(1,\log t,\ldots,(\log t)^{r-1}\), and the second from
\(1,\log(1-t),\ldots,(\log(1-t))^{r-1}\). Hence
\begin{equation}\label{cg:log-vandermonde-integral}
\begin{split}
 \det M
 ={}&\int_{\Omega_r}
   \left(\prod_{k=1}^{r}w_{x,y}(t_k)\right)
   \V(\log t_1,\ldots,\log t_r)\\
 &\hspace{24mm}\cdot
   \V(\log(1-t_1),\ldots,\log(1-t_r))
   \dd\mathbf t .
\end{split}
\end{equation}
Absolute integrability follows by expanding the two Vandermonde
determinants: every term is a product of the finite absolute logarithmic
moments appearing in \eqref{cg:beta-derivatives}.
On \(\Omega_r\), the values \(\log t_k\) are strictly increasing, whereas
\(\log(1-t_k)\) are strictly decreasing. Thus the first Vandermonde is
positive and the second has the strict sign \(\ep_r\). The positive
weight and the nonempty open integration domain prove
\begin{equation}\label{cg:strict-beta-derivative}
 \ep_r\det M>0.
\end{equation}

It remains to transfer this determinant to derivatives of \(f\).
For \(0\le i,p\le r-1\), define the lower triangular matrix \(L_x\) by
\[
 (L_x)_{ip}=
 \begin{cases}
  \displaystyle\binom{i}{p}\Gamma^{(i-p)}(x),&p\le i,\\
  0,&p>i,
 \end{cases}
\]
and define \(L_y\) analogously. Let
\(\mathcal H=(f^{(p+q)}(x+y))_{p,q=0}^{r-1}\).
Since \(B(x,y)=\Gamma(x)\Gamma(y)f(x+y)\), the two-variable Leibniz
rule is exactly the matrix identity
\begin{equation}\label{cg:leibniz}
 M=L_x\mathcal H L_y^{\mathsf T}.
\end{equation}
Indeed, its \((i,j)\)-entry is
\[
 \sum_{p=0}^{i}\sum_{q=0}^{j}
 \binom{i}{p}\binom{j}{q}
 \Gamma^{(i-p)}(x)\Gamma^{(j-q)}(y)
 f^{(p+q)}(x+y),
\]
which is \(\partial_x^i\partial_y^j B(x,y)\).
Every diagonal entry of \(L_x\) equals \(\Gamma(x)\), and every
diagonal entry of \(L_y\) equals \(\Gamma(y)\). Consequently,
\begin{equation}\label{cg:leibniz-determinant}
 \det M=[\Gamma(x)\Gamma(y)]^r H_r(x+y).
\end{equation}
Combining \eqref{cg:log-vandermonde-integral} with
\eqref{cg:leibniz-determinant} proves
\eqref{cg:confluent-integral}. The scaling factor is strictly positive,
so \eqref{cg:strict-beta-derivative} proves
\eqref{cg:confluent-inequality}.
\end{proof}

For example, the \(r=2\) conclusion is the strict differential
inequality
\[
 f(s)f''(s)-[f'(s)]^2<0,
\]
and the theorem supplies its complete higher-order determinant hierarchy.

\Needspace{11\baselineskip}
\subsection{Inverse factorials of arbitrary real order}

\begin{corollary}[Real-order inverse factorials]\label{cg:inverse-factorial}
For \(x>0\) and \(x+u>0\), define
\[
 (x)_u=\frac{\Gamma(x+u)}{\Gamma(x)}.
\]
If \(0<x_1<\cdots<x_r\) and \(u_1<\cdots<u_r\) are real, with
\(x_1+u_1>0\), then
\begin{equation}\label{cg:inverse-factorial-inequality}
 \ep_r\det\left[\frac1{(x_i)_{u_j}}\right]_{i,j=1}^{r}>0.
\end{equation}
In particular, the inverse-factorial kernel is \(\SSR_\infty\) on
\((0,\infty)\times[0,\infty)\), with no integrality assumption on the
order \(u\).
\end{corollary}

\begin{proof}
Each row is a positive multiple of a reciprocal-gamma row:
\[
 \det\left[\frac1{(x_i)_{u_j}}\right]
 =
 \left(\prod_{i=1}^{r}\Gamma(x_i)\right)
 \det\left[\frac1{\Gamma(x_i+u_j)}\right].
\]
The factor is positive, and Theorem~\ref{cg:main} applies because
\(x_1+u_1>0\).
\end{proof}

\section{Gamma-ratio kernels: complete global classification}\label{gr:section}


Karp, Vishnyakova, and Zhang proposed all-order sign regularity of the
gamma-ratio kernel for \(c>d\) in Section~3, Example~11 of
\cite{KVZ}. The same example proposed the order-three signature
\((+,-,-)\). We prove the all-order assertion on the natural
positive-argument domain and determine its parameter-dependent signature.
The proposed third sign requires correction when \(0<c-d<1\).

\begin{theorem}[Complete gamma-ratio classification]\label{thm:main}
Let \(c,d\in\R\), set \(\delta=c-d\), and let
\(x_1<\cdots<x_r\) and \(y_1<\cdots<y_r\) be real nodes satisfying
\begin{equation}\label{eq:domain}
 x_1+y_1+\min(c,d)>0.
\end{equation}
Define
\[
 D_r=\det\left[
 \frac{\Gamma(c+x_i+y_j)}{\Gamma(d+x_i+y_j)}
 \right]_{i,j=1}^{r}.
\]
Then the following statements hold.
\begin{enumerate}[label=\textnormal{(\roman*)}]
\item If \(\delta\notin\Z_{\ge0}\), then
\begin{equation}\label{eq:classification}
 \sgn D_r=\sgn C_r(\delta),\qquad D_r\ne0.
\end{equation}
In particular, if \(\delta<0\), then \(D_r>0\) for every \(r\).
\item If \(\delta=m\in\Z_{\ge0}\), then
\begin{equation}\label{eq:integer}
 \ep_r D_r>0\quad(1\le r\le m+1),
 \qquad D_r=0\quad(r\ge m+2).
\end{equation}
On every nonempty open rectangle satisfying the positive-argument
condition, the kernel has exact rank \(m+1\).
\end{enumerate}
\end{theorem}

Condition~\eqref{eq:domain} means that both gamma arguments are positive
for every entry. No integrality or equal-spacing assumption is imposed on
either node list. In particular, the theorem applies to every
\(c>d>0\) and all increasing positive real nodes.

For noninteger \(\delta>0\), set \(m=\lfloor\delta\rfloor\). The signature
can equivalently be written
\begin{equation}\label{eq:floor-sign}
 \sigma_r(\delta)=
 (-1)^{\sum_{j=0}^{r-1}\min(j,m+1)}.
\end{equation}
Indeed, the factor \((-\delta)_j\) has exactly
\(\min(j,m+1)\) negative terms. Some initial signatures are displayed in
Table~\ref{tab:signs}. A zero in the table means that every minor of that
order vanishes.

\begin{table}[htbp]
\centering
\caption{The signs of all ordered minors, for the first six orders.}
\label{tab:signs}
\begin{tabular}{lcccccc}
\toprule
Parameter difference & \(r=1\)&\(r=2\)&\(r=3\)&\(r=4\)&\(r=5\)&\(r=6\)\\
\midrule
\(\delta<0\)   & \(+\)&\(+\)&\(+\)&\(+\)&\(+\)&\(+\)\\
\(\delta=0\)   & \(+\)&\(0\)&\(0\)&\(0\)&\(0\)&\(0\)\\
\(0<\delta<1\) & \(+\)&\(-\)&\(+\)&\(-\)&\(+\)&\(-\)\\
\(\delta=1\)   & \(+\)&\(-\)&\(0\)&\(0\)&\(0\)&\(0\)\\
\(1<\delta<2\) & \(+\)&\(-\)&\(-\)&\(-\)&\(-\)&\(-\)\\
\(\delta=2\)   & \(+\)&\(-\)&\(-\)&\(0\)&\(0\)&\(0\)\\
\(2<\delta<3\) & \(+\)&\(-\)&\(-\)&\(+\)&\(-\)&\(+\)\\
\(\delta=3\)   & \(+\)&\(-\)&\(-\)&\(+\)&\(0\)&\(0\)\\
\(3<\delta<4\) & \(+\)&\(-\)&\(-\)&\(+\)&\(+\)&\(+\)\\
\bottomrule
\end{tabular}
\end{table}



\subsection{Gamma densities and preservation of determinant signs}

For \(a,u>0\), define the gamma density
\begin{equation}\label{eq:density}
 p_a(u)=\frac{u^{a-1}e^{-u}}{\Gamma(a)}.
\end{equation}
The normalization follows from Euler's gamma integral; the beta and gamma
integrals used below are recorded in \cite{DLMF}.

\begin{lemma}\label{lem:density}
The kernel \((a,u)\mapsto p_a(u)\) is \(\STP_\infty\) on
\((0,\infty)^2\).
\end{lemma}
\begin{proof}
Its ordered matrices are positive row and column scalings of
\([e^{a_i\log u_j}]\). Their determinants are strictly positive by
Lemma~\ref{lem:monomials}; transposition does not change a determinant.
\end{proof}

\begin{lemma}[The gamma moment factorization]\label{lem:moment}
If \(a,b>0\) and \(a+b+\delta>0\), then
\begin{equation}\label{eq:mixture}
 \frac{\Gamma(a+b+\delta)}{\Gamma(a+b)}
 =\int_0^\infty\!\int_0^\infty
 p_a(u)(u+v)^\delta p_b(v)\dd v\dd u.
\end{equation}
The integral is finite and positive.
\end{lemma}
\begin{proof}
Set \(s=u+v\) and \(t=u/(u+v)\). Then
\(u=st\), \(v=s(1-t)\), and the Jacobian is \(s\). The integral is
\[
 \frac{1}{\Gamma(a)\Gamma(b)}
 \left(\int_0^\infty e^{-s}s^{a+b+\delta-1}\dd s\right)
 \left(\int_0^1 t^{a-1}(1-t)^{b-1}\dd t\right).
\]
The first integral is \(\Gamma(a+b+\delta)\), and the second is
\(\Gamma(a)\Gamma(b)/\Gamma(a+b)\). The stated hypotheses are
exactly the needed endpoint conditions for these two positive integrals.
\end{proof}

The following composition formula is the mechanism that transfers the
power-kernel signs. Write
\(\simplex r=\{(u_1,\ldots,u_r):0<u_1<\cdots<u_r<\infty\}\).

\begin{proposition}[Double gamma-density composition]\label{prop:composition}
Let \(H:(0,\infty)^2\to\R\) be continuous and assume
\[
 \int_0^\infty\!\int_0^\infty
 p_a(u)|H(u,v)|p_b(v)\dd v\dd u<\infty
\]
for the shape pairs under consideration. Put
\(T_H(a,b)=\int\!\int p_a(u)H(u,v)p_b(v)\dd v\dd u\).
For increasing shape lists \(a_i,b_j>0\), if the above integrability
holds for every pairing \((a_i,b_j)\), then
\begin{align}\label{eq:double-composition}
 \det[T_H(a_i,b_j)]
 &=\int_{\simplex r}\!\int_{\simplex r}
 \det[p_{a_i}(u_k)]_{i,k=1}^{r}\,
 \det[H(u_k,v_l)]_{k,l=1}^{r}\notag\\[-0.2em]
 &\hspace{27mm}\cdot
 \det[p_{b_j}(v_l)]_{l,j=1}^{r}
 \dd\mathbf v\dd\mathbf u.
\end{align}
If every ordered \(r\)-minor of \(H\) has one strict sign, then the
left side has the same strict sign. The corresponding non-strict
preservation statement also holds.
\end{proposition}

\begin{proof}
The usual continuous Cauchy--Binet identity is obtained by expanding two
finite determinants, applying Fubini, and relabeling integration
variables. Applying it twice gives the full product-space version of
\eqref{eq:double-composition}, with factor \(1/(r!)^2\).
The three-determinant integrand is symmetric separately in the \(u\)'s
and the \(v\)'s: each permutation changes the signs of two determinants.
Restricting both integrations to their increasing simplices therefore
cancels the factorial factor.

For completeness, absolute integrability can be checked before making
these expansions. The absolute value of the product of three
determinants is bounded by a finite sum of at most \((r!)^3\) terms.
Each term is a product of gamma densities and factors
\(|H(u_k,v_{\rho(k)})|\), where \(\rho\) is a permutation. Pairing
\(u_k\) with \(v_{\rho(k)}\), its full-space integral factors as a
product of \(r\) finite integrals from the hypothesis. Thus every
interchange is justified.

By Lemma~\ref{lem:density}, the two outside determinants in
\eqref{eq:double-composition} are strictly positive. If the middle
determinant has a fixed strict sign, the signed integrand is continuous
and strictly positive everywhere on the nonempty open product of
simplices. Its integral is therefore strictly positive. This also proves
the non-strict assertion when strict positivity is relaxed.
\end{proof}

\subsection{Proof of the global classification}

We now prove Theorem~\ref{thm:main}. First take increasing positive shape
lists \(a_i,b_j\) with
\(a_1+b_1+\delta>0\). By Lemma~\ref{lem:moment}, the gamma-ratio
matrix is the double transform of \(P_\delta(u,v)=(u+v)^\delta\).
Every pairing satisfies the integrability hypotheses of
Proposition~\ref{prop:composition}. Thus
\begin{align}\label{eq:gamma-determinant-integral}
 &\det\left[
 \frac{\Gamma(a_i+b_j+\delta)}{\Gamma(a_i+b_j)}
 \right]_{i,j=1}^{r}\notag\\
 &\quad=\int_{\simplex r}\!\int_{\simplex r}
 \det[p_{a_i}(u_k)]
 \det[(u_k+v_l)^\delta]
 \det[p_{b_j}(v_l)]
 \dd\mathbf v\dd\mathbf u.
\end{align}

If \(\delta\notin\Z_{\ge0}\), Proposition~\ref{prop:power} gives
the strict middle sign \(\sgn C_r(\delta)\). The other two determinants
are positive, so \eqref{eq:classification} follows, including strictness.
For \(\delta<0\), every rising factorial \((-\delta)_j\) is positive,
which proves strict total positivity.

For \(\delta=m\in\Z_{\ge0}\), the same argument proves strict sign
\(\ep_r\) through order \(m+1\). In addition,
\[
 \frac{\Gamma(s+m)}{\Gamma(s)}=(s)_m
\]
is a polynomial of degree \(m\). A polynomial in \(a+b\) of degree
\(m\) is a sum of \(m+1\) functions of \(a\) times functions of \(b\).
Hence every higher minor vanishes. A strictly nonzero minor of order
\(m+1\) proves that the rank is exactly \(m+1\).

It remains to pass to arbitrary real node lists. By
\eqref{eq:domain}, choose a real \(h\) with
\[
 -x_1<h<d+y_1.
\]
Set \(a_i=x_i+h\) and \(b_j=y_j+d-h\). Both lists are positive and
increasing, and
\[
 a_i+b_j=d+x_i+y_j>0,\qquad
 a_i+b_j+\delta=c+x_i+y_j>0.
\]
The already proved shape-variable result applies to precisely the
requested entries. This proves the full theorem.\hfill\(\square\)

\begin{remark}[Why no strictness is lost]
The argument does not infer strict signs from a limit of non-strict
inequalities. Strictness holds pointwise in the determinant integrand
on an open set of positive measure. This is the decisive distinction
between merely obtaining \(\sigma_rD_r\ge0\) and proving
\(\sigma_rD_r>0\).
\end{remark}

\subsection{Exact formulas and the correction at order three}

\begin{proposition}[Consecutive columns]\label{prop:consecutive}
Suppose \(y_j=y_0+j\), \(j=0,\ldots,r-1\), and all gamma arguments
are positive. For arbitrary increasing real row nodes,
\begin{align}\label{eq:consecutive}
 &\det_{\substack{1\le i\le r\\0\le j\le r-1}}
 \left[\frac{\Gamma(x_i+c+y_0+j)}
 {\Gamma(x_i+d+y_0+j)}\right]\notag\\
 &\qquad=
 \V(x_1,\ldots,x_r)
 \prod_{i=1}^{r}
 \frac{\Gamma(x_i+c+y_0)}{\Gamma(x_i+d+y_0+r-1)}
 \prod_{j=0}^{r-1}(d-c)_j.
\end{align}
\end{proposition}

\begin{proof}
Put \(t_i=x_i+y_0\). Factor the positive quantity
\(\Gamma(t_i+c)/\Gamma(t_i+d+r-1)\) from row \(i\).
The remaining entries are evaluations of the polynomials
\[
 p_j(t)=(t+c)_j(t+d+j)_{r-1-j},\qquad 0\le j\le r-1.
\]
Their evaluation determinant is \(\V(t)A_r(c,d)\), where \(A_r\)
is the coefficient determinant in increasing monomial order. Translation
to powers of \(t+c\) has determinant one. At \(t=-c\),
\(p_0(-c)=(d-c)_{r-1}\), while \(p_j(-c)=0\) for \(j\ge1\).
Expand the coefficient determinant along its constant row and divide
the remaining polynomials by \(t+c\). The resulting family is the
family of size \(r-1\) with parameters \(c+1,d+1\). Consequently,
\[
 A_r(c,d)=(d-c)_{r-1}A_{r-1}(c+1,d+1).
\]
Starting from \(A_1=1\) proves
\(A_r=\prod_{j=0}^{r-1}(d-c)_j\). Finally \(\V(t)=\V(x)\).
\end{proof}

The consecutive-column identity first determines a necessary signature.
Theorem~\ref{thm:main} is the step that proves the same signature for
arbitrary columns as well.

At order three,
\begin{equation}\label{eq:third}
 C_3(\delta)=\delta^2(1-\delta).
\end{equation}
Thus \(0<\delta<1\) gives positive third-order minors, whereas
\(\delta>1\) gives negative third-order minors. For \(\delta=1\), all
third-order minors vanish. These conclusions now hold for every allowed
choice of nodes, not only for consecutive columns.

\begin{example}[An exact positive third-order determinant]
Take \(c=1\), \(d=1/2\), and \(x_i=y_i=1/4+i\), \(i=0,1,2\).
Then the matrix is
\[
 \frac{\sqrt\pi}{2}
 \begin{pmatrix}
 1&3/2&15/8\\
 3/2&15/8&35/16\\
 15/8&35/16&315/128
 \end{pmatrix},
\]
and its determinant equals
\[
 \frac{\pi^{3/2}}{8}\cdot\frac{5}{1024}
 =\frac{5\pi^{3/2}}{8192}>0.
\]
This disproves the parameter-independent negative third sign proposed
in \cite[Section~3, Example~11]{KVZ}. Theorem~\ref{thm:main} nevertheless
establishes the all-order sign-regularity assertion with the corrected
signature.
\end{example}

\section{Failure of pointwise product closure}\label{pc:section}

The integral composition above preserves determinant signs under its
explicit hypotheses. Pointwise multiplication has a different behavior.
The general product expectation in \cite[Section~3, Example~11]{KVZ}
fails even for strictly sign-regular factors.

\begin{theorem}[A product with second-order minors of both signs]
\label{pc:counterexample}
For every \(0<q<1\) and \(a>0\), the kernels on \((0,\infty)^2\)
\[
 K_1(x,y)=e^{a(x+y)^2},\qquad
 K_2(x,y)=\frac1{\Gamma_q(x+y)}
\]
are respectively \(\STP_\infty\) and \(\SSR_\infty\).
Nevertheless their pointwise product is not sign regular of order two:
it has positive and negative ordered second-order minors.
\end{theorem}

\begin{proof}
Factor
\[
 K_1(x,y)=e^{ax^2}e^{2axy}e^{ay^2}.
\]
The exponential kernel is strictly totally positive by
Lemma~\ref{lem:monomials}, after setting \(t=e^{2ax}\).
The outside factors are positive row and column multipliers.
Theorem~\ref{qg:main} gives the claimed strict signature of \(K_2\).

Set \(h(s)=e^{as^2}/\Gamma_q(s)\) and \(E=e^{2a}>1\).
For the two identical node lists
\((s/2,s/2+1)\), where \(s>0\), the product kernel has determinant
\[
 D(s)=h(s)h(s+2)-h(s+1)^2.
\]
The recurrence \eqref{qg:recurrence} gives
\[
 \frac{h(s)h(s+2)}{h(s+1)^2}
 =E\frac{1-q^s}{1-q^{s+1}},
\]
and therefore the exact formula
\begin{equation}\label{pc:determinant}
 D(s)=h(s+1)^2
 \frac{(E-1)-(E-q)q^s}{1-q^{s+1}}.
\end{equation}
The prefactor and denominator are positive. Since
\(0<(E-1)/(E-q)<1\), the number
\begin{equation}\label{pc:threshold}
 s_* =\frac{\log((E-1)/(E-q))}{\log q}
\end{equation}
is positive. As \(q^s\) strictly decreases, \eqref{pc:determinant}
is negative for \(0<s<s_*\), zero at \(s_*\), and positive for
\(s>s_*\). These are admissible ordered minors of the same kernel.
No single second-order signature is possible.
\end{proof}

\begin{example}[Two exact minors of one product]
Take \(q=1/2\) and \(a=(\log2)/2\), so that \(E=2\).
Then
\[
 D(1/2)=h(3/2)^2\frac{2\sqrt2-3}{2\sqrt2-1}<0.
\]
On the other hand, \(h(1)=\sqrt2\), \(h(2)=4\), and
\(h(3)=32\sqrt2/3\), hence
\[
 D(1)=\frac{16}{3}>0.
\]
The respective node lists are \((1/4,5/4)\) and \((1/2,3/2)\).
Thus the sign change is certified by exact values, without a numerical
search or an asymptotic approximation.
\end{example}

This counterexample concerns multiplication of sign-regular kernels
with potentially different signatures. It makes no claim that the
product of two strictly totally positive kernels fails total
positivity; the second factor above already has negative second-order
minors. The gamma-ratio proof uses integral composition and is independent
of a pointwise product-closure assertion.

\section{Zeros, interpolation, and ratios}\label{app:section}

\subsection{A common consequence of strict sign regularity}

Let \(I,J\) be nonempty real intervals, with \(I\) open, and let
\(K:I\times J\to(0,\infty)\) be a continuous strictly sign-regular
kernel of every order. For increasing shifts \(\lambda_j\in J\), put
\(k_j(x)=K(x,\lambda_j)\). For a finite or infinite real coefficient
sequence \(\mathbf a\), let \(S^-(\mathbf a)\) be the number of sign
changes after zero entries are deleted. In the infinite case the next
result requires this number to be finite.

\begin{theorem}[A generalized Descartes rule]\label{cor:zeros}
Let \(\mathbf a\) be nonzero and \(s=S^-(\mathbf a)<\infty\).
Suppose
\begin{equation}\label{app:abs-convergence}
 \sum_j |a_j|K(x,\lambda_j)<\infty\qquad(x\in I).
\end{equation}
Then \(F(x)=\sum_j a_jK(x,\lambda_j)\) is not identically zero and
has at most \(s\) distinct zeros in \(I\).
For a finite sum of \(N\) terms, this gives at most \(N-1\) zeros.
\end{theorem}

\begin{proof}
Put a separator \(\theta_l\in J\) between the consecutive nonzero
coefficient blocks at the \(l\)-th sign change, for \(1\le l\le s\).
Choose each separator strictly between the adjacent active shifts, so
none coincides with an active shift. Suppose \(F\) has distinct zeros
\(z_0<\cdots<z_s\). Form
\begin{equation}\label{app:separator}
 Q(t)=\det\begin{bmatrix}
 K(z_0,\theta_1)&\cdots&K(z_0,\theta_s)&K(z_0,t)\\
 \vdots&&\vdots&\vdots\\
 K(z_s,\theta_1)&\cdots&K(z_s,\theta_s)&K(z_s,t)
 \end{bmatrix}.
\end{equation}
By strict sign regularity, \(Q(t)\ne0\) away from the separators.
Sorting its last column into increasing order changes the permutation
sign once whenever \(t\) crosses a separator. Thus all the nonzero
numbers \(a_jQ(\lambda_j)\) have the same strict sign. Their sum is
nonzero. Expansion of \(Q\) along its last column instead gives
\[
 \sum_j a_jQ(\lambda_j)
 =\sum_{i=0}^s A_iF(z_i)=0
\]
for fixed cofactors \(A_i\), a contradiction. This expansion is valid
also for an infinite sum: \eqref{app:abs-convergence} at the finitely
many \(z_i\) implies
\(\sum_j|a_jQ(\lambda_j)|
 \le\sum_i|A_i|\sum_j|a_j|K(z_i,\lambda_j)<\infty\).
The same argument covers \(s=0\). An identically zero function would
also supply the prohibited \(s+1\) zeros.
\end{proof}

The theorem applies to both reciprocal kernels by
Theorems~\ref{qg:main} and~\ref{cg:main}, and to the gamma-ratio
kernel whenever \(c-d\notin\Z_{\ge0}\), by Theorem~\ref{thm:main}.
Choose \(I\) and \(J\) so that the corresponding positive-argument
condition holds throughout their product. The proof needs strict
minors only through order \(s+1\).

\begin{corollary}[Unique finite interpolation]\label{cor:interpolation}
For any \(N\) distinct points in \(I\) and any prescribed real
values at those points, there is exactly one coefficient vector
\((a_1,\ldots,a_N)\) whose function \(\sum_{j=1}^Na_jk_j\) takes
those values.
\end{corollary}
\begin{proof}
The interpolation matrix is a strictly nonzero ordered minor of \(K\).
\end{proof}

This is an assertion of exact invertibility, not a uniform bound on
numerical condition numbers. The zero count concerns distinct zeros.
For the integer gamma-ratio case \(c-d=m\ge0\), the span of \(N\)
distinct shifts instead has dimension \(\min(N,m+1)\): the strict
minors through order \(m+1\) give the lower bound, and the polynomial
representation gives the upper bound. Every nonzero member has degree
at most \(m\), hence at most \(m\) distinct real zeros. Unique
\(N\)-point interpolation still holds for \(N\le m+1\), whereas
coefficient injectivity fails for \(N>m+1\).

\subsection{Convergent reciprocal \texorpdfstring{$q$}{q}-gamma series}

For \(K(x,y)=1/\Gamma_q(x+y)\), a useful sufficient condition for
\eqref{app:abs-convergence}, locally uniformly on the positive-sum
domain, is
\begin{equation}\label{app:q-convergence}
 \sum_j |a_j|(1-q)^{\lambda_j}<\infty.
\end{equation}
Indeed, the product formula gives
\[
 0<K(x,\lambda_j)
 \le\frac{(1-q)^{x-1}}{(q;q)_\infty}(1-q)^{\lambda_j}.
\]
On every compact \(x\)-interval the factor preceding
\((1-q)^{\lambda_j}\) is bounded. The Weierstrass test proves the
claimed uniform absolute convergence. Thus the zero theorem applies
to infinite reciprocal \(q\)-gamma series under
\eqref{app:q-convergence} and a finite number of coefficient sign
changes. It makes no general injectivity claim for an infinite
coefficient sequence with infinitely many sign changes.

\subsection{Level crossings and monotone coefficient ratios}

For the finite family above, take \(b_j>0\) and define
\begin{equation}\label{eq:ratio}
 R(x)=\frac{A(x)}{B(x)},\qquad
 A(x)=\sum_{j=1}^N a_jk_j(x),\quad
 B(x)=\sum_{j=1}^N b_jk_j(x)>0,\quad
 r_j=\frac{a_j}{b_j}.
\end{equation}

\begin{corollary}[Level crossings and monotonicity]\label{cor:ratios}
Under the all-order strictness hypotheses of this section, for any
real \(t\), unless all \(r_j=t\),
\begin{equation}\label{app:level-bound}
 \#\{x\in I:R(x)=t\}\le S^-((r_j-t)_{j=1}^N).
\end{equation}
If only strict sign regularity through order two is assumed, with
second sign \(\sigma_2\), a nonconstant nondecreasing sequence
\((r_j)\) makes \(R\) strictly increasing when \(\sigma_2=+1\)
and strictly decreasing when \(\sigma_2=-1\). The directions reverse
for a nonconstant nonincreasing sequence.
\end{corollary}

\begin{proof}
Apply Theorem~\ref{cor:zeros} to
\(A-tB=\sum_j b_j(r_j-t)k_j\). The positive denominator identifies
its zeros with the requested level crossings.
For the second assertion, direct pairing of terms gives, for
\(x_1<x_2\),
\begin{align}\label{app:pair-identity}
 &B(x_1)B(x_2)\{R(x_2)-R(x_1)\}\notag\\
 &\qquad=\sum_{i<j}b_ib_j(r_j-r_i)
 \det\begin{pmatrix}k_i(x_1)&k_j(x_1)\\
 k_i(x_2)&k_j(x_2)\end{pmatrix}.
\end{align}
Every determinant has strict sign \(\sigma_2\). Monotonicity of
\((r_j)\) makes all summands have one weak sign, and nonconstancy
makes at least one of them strict.
\end{proof}

For either reciprocal kernel, \(\sigma_2=-1\); increasing
coefficient ratios therefore give a decreasing function ratio.
For gamma ratios, \(\sigma_2=\sgn(d-c)\) whenever \(c\ne d\).
This includes positive integer differences, since the second-order
minors are still strict. At \(c=d\), the kernel is constant.

The level-crossing assertion also applies to an infinite ratio when
the numerator and denominator converge absolutely, the denominator
has strictly positive coefficients, and
\(S^-((a_j/b_j-t)_j)<\infty\). The monotonicity assertion extends
as well when both coefficient series converge absolutely at every
point: the paired double sum in \eqref{app:pair-identity} is then
absolutely convergent, and any strict pair supplies a strict term.
For the reciprocal \(q\)-gamma case, condition
\eqref{app:q-convergence} for both \((a_j)\) and \((b_j)\) is
sufficient.

\subsection{A probability interpretation}

For a finite positive denominator, define
\[
 p_j(x)=\frac{b_jK(x,\lambda_j)}{B(x)},\qquad
 \sum_jp_j(x)=1.
\]
Then \(R(x)=\sum_jp_j(x)r_j\) is the expectation of the
coefficient ratios. For \(x_1<x_2\) and \(i<j\), division of an
ordered second minor by positive entries shows that
\[
 \sgn\left(
 \frac{K(x_2,\lambda_j)}{K(x_1,\lambda_j)}
 -\frac{K(x_2,\lambda_i)}{K(x_1,\lambda_i)}\right)=\sigma_2.
\]
The likelihood ratios \(p_j(x_2)/p_j(x_1)\) therefore increase
with \(j\) for \(\sigma_2=+1\) and decrease for \(\sigma_2=-1\).
Equation~\eqref{app:pair-identity} is an exact ordering result for
expectations in this family. The determinant theorem provides this
structural control without asserting numerical stability or an
additional statistical estimation theorem.

\subsection{The sign pattern when the zero bound is attained}

\begin{lemma}[Oriented maximal-zero pattern]\label{app:orientation}
Under the hypotheses of Theorem~\ref{cor:zeros}, suppose \(s\ge1\)
and \(F\) has \(s\) distinct zeros \(z_1<\cdots<z_s\). Let
\(\alpha\) be the sign of the first nonzero coefficient, and let
\(\sigma_r\) be the kernel signature. For \(x\) distinct from these
zeros, put \(l=\#\{i:z_i<x\}\). Then
\begin{equation}\label{app:oriented-sign}
 \sgn F(x)=\alpha\sigma_s\sigma_{s+1}(-1)^l.
\end{equation}
Only strict sign regularity through order \(s+1\) is required for
this statement and its accompanying zero bound.
\end{lemma}

\begin{proof}
Choose separators \(\theta_1<\cdots<\theta_s\) strictly between
the successive coefficient-sign blocks, as in the preceding proof.
Form \(Q_x(t)\) from rows \((z_1,\ldots,z_s,x)\) and columns
\((\theta_1,\ldots,\theta_s,t)\) of \(K\). If
\(k=\#\{j:\theta_j<t\}\), sorting these rows and columns gives
\[
 \sgn Q_x(t)=\sigma_{s+1}(-1)^{l+k}.
\]
For an active coefficient at \(t=\lambda_j\), its sign is
\(\alpha(-1)^k\). Consequently the nonzero summands of
\(\sum_j a_jQ_x(\lambda_j)\) all have sign
\(\alpha\sigma_{s+1}(-1)^l\). Expansion along the last column,
using \(F(z_i)=0\), gives
\[
 \sum_j a_jQ_x(\lambda_j)
 =F(x)\det[K(z_i,\theta_j)]_{i,j=1}^s.
\]
The latter determinant has sign \(\sigma_s\), proving the formula.
The same finite-cofactor bound used in Theorem~\ref{cor:zeros}
justifies every sum in the infinite case.
\end{proof}

\subsection{Unimodality and its parameter transition}

A \emph{peak} is a nondecreasing part followed by a nonincreasing
part; a \emph{trough} has the reversed order. Monotone and constant
functions or sequences are admitted as degenerate cases. Equivalently,
a continuous peak-shaped function on an interval has no valley triple:
there are no \(x_1<x_2<x_3\) with
\(R(x_2)<\min(R(x_1),R(x_3))\).

\begin{theorem}[Finite unimodality transfer]\label{app:unimodality}
Let \(K\) be a positive continuous \(\SSR_3\) kernel, and let
\(R\) be the finite ratio \eqref{eq:ratio}. If the coefficient ratios
\((r_j)\) have one peak or one trough, then \(R\) is unimodal.
Put \(\eta=\sigma_2\sigma_3\). If \(R\) is not monotone, its
pattern agrees with the coefficient pattern when \(\eta=+1\), and
is reversed when \(\eta=-1\).
\end{theorem}

\begin{proof}
Suppose first that \((r_j)\) is peak-shaped. For every level \(t\),
the coefficients \(b_j(r_j-t)\) have at most two sign changes.
If they have two, their signs occur in the pattern \((-,+,-)\).
The proof of Theorem~\ref{cor:zeros}, using only minors through
order three, bounds the zeros of \(R-t\) by that number, unless
all \(r_j=t\), in which case \(R\) is constant.

When two zeros occur, Lemma~\ref{app:orientation} with \(s=2\)
forces the sign pattern of \(R-t\) to be
\(\eta(-,+,-)\), since its denominator is positive.
If \(\eta=+1\), a valley triple would allow a level strictly
between its middle value and its two endpoint values. Continuity
would then give two zeros with pattern \((+,-,+)\), a contradiction.
Thus \(R\) has no valley triple and is peak-shaped. If \(\eta=-1\),
the same argument excludes peak triples and makes \(R\) trough-shaped.
Replacing \((a_j)\) by \((-a_j)\) proves the assertion for a
trough-shaped coefficient sequence.
\end{proof}

This is the finite-sum conclusion of \cite[Theorem~2.5]{KVZ}, with
a direct proof included here. For infinite sums the same argument
requires continuity of the ratio as well as absolute convergence;
local uniform convergence and a positive denominator suffice.

\begin{corollary}[The signatures determine the shape]
\label{cor:unimodal}
For the reciprocal gamma and reciprocal \(q\)-gamma kernels, a
single coefficient peak or trough is preserved whenever the resulting
ratio is not monotone. For the gamma-ratio kernel, write \(\delta=c-d\).
Its nonmonotone patterns obey the following rule:
\begin{enumerate}[label=\textnormal{(\roman*)}]
\item If \(\delta<0\), the coefficient pattern is preserved.
\item If \(0<\delta<1\), a peak becomes a trough and a trough
becomes a peak.
\item If \(\delta>1\), the coefficient pattern is preserved.
\end{enumerate}
At \(\delta=0\), every such ratio is constant. At \(\delta=1\),
it is monotone or constant.
\end{corollary}

\begin{proof}
Both reciprocal kernels have \(\sigma_2=\sigma_3=-1\), hence
\(\eta=+1\). For gamma ratios with \(\delta<0\), both signs are
positive. For \(\delta>0\), \(\delta\ne1\), Theorem~\ref{thm:main}
gives \(\sigma_2=-1\) and \(\sigma_3=\sgn(1-\delta)\).
Theorem~\ref{app:unimodality} proves the stated regimes, including
integer \(\delta\ge2\), where the first three orders remain strict.
If \(\delta=0\), every kernel entry is one. If \(\delta=1\),
\(k_j(x)=x+\lambda_j+d\); the ratio is a quotient of two affine
functions with positive denominator, whose derivative has constant
sign or vanishes identically.
\end{proof}

\begin{example}[An exact peak-to-trough reversal]\label{ex:trough}
Let \(c=1\), \(d=1/2\), take shifts \((0,1,2)\), and set
\[
 (a_1,a_2,a_3)=(0,1,0),\qquad
 (b_1,b_2,b_3)=(1,1,2).
\]
The coefficient ratios are \((0,1,0)\), a peak. For \(x>0\), put
\(s=x+1/2>1/2\). Gamma recurrence cancels a common positive factor
in \eqref{eq:ratio} and yields exactly
\begin{equation}\label{eq:trough-formula}
 R(x)=\frac{2s^2+3s+1}{8s^2+13s+4},\qquad
 R'(x)=\frac{2s^2-1}{(8s^2+13s+4)^2}.
\end{equation}
Thus \(R\) strictly decreases on
\(0<x<(\sqrt2-1)/2\) and strictly increases after that point.
It has one strict interior trough at
\[
 x_*=(\sqrt2-1)/2.
\]
This example shows that the change in the third determinant sign has
an observable consequence for the shape of function ratios. Retaining
the proposed universal signature \((+,-,-)\) would give the wrong
pattern in this regime.
\end{example}

\section{A nonlinear question for reciprocal \texorpdfstring{\(q\)}{q}-gamma series}
\label{ob:nonlinear}

Sign regularity controls linear combinations, their zeros, and the level
crossings of ratios. Logarithmic concavity of a sum is a different
assertion: it compares products of sums at different parameter values.
We identify an exact additional inequality that would establish the
positive-parameter version for reciprocal \(q\)-gamma series. The
identity leading to that inequality is proved below; the inequality
itself is not proved in this paper.

\subsection{Admissible coefficients and convergence}
\label{ob:admissible}

Fix \(0<q<1\). Call a sequence
\((g_n)_{n\ge0}\) \emph{doubly positive} if it is nonzero,
nonnegative, log-concave,
\[
 g_n^2\ge g_{n-1}g_{n+1}\qquad(n\ge1),
\]
and has no internal zeros. The last condition means that its positive
support is an interval of nonnegative integers, possibly finite. Initial or terminal
zeros are permitted. For the proposed nonlinear inequality we require
such a sequence and choose \(z\ge0\) with
\begin{equation}\label{ob:mass}
 0<M_g(z):=\sum_{n=0}^{\infty}g_n\bigl((1-q)z\bigr)^n<\infty.
\end{equation}
At \(z=0\), this means \(g_0>0\); if \(g_0=0\), the resulting
series is identically zero and its Tur\'anian vanishes. A finite
coefficient sequence satisfies the finiteness requirement for every
finite \(z\). More generally, if \(\rho\) is the radius of
convergence of \(\sum_n g_n w^n\), it suffices to take
\(0<z<\rho/(1-q)\); condition~\eqref{ob:mass} also permits a
convergent boundary point.

Write
\begin{equation}\label{ob:series}
 G(\mu;z)=\sum_{n=0}^{\infty}g_nz^n\kappa_q(\mu+n),
 \qquad \kappa_q(t)=\frac1{\Gamma_q(t)},\qquad \mu>0.
\end{equation}
The product representation \cite[Section~5.18]{DLMF} is
\begin{equation}\label{ob:reciprocal-product}
 \kappa_q(t)=\frac{(1-q)^{t-1}(q^t;q)_\infty}{(q;q)_\infty},
 \qquad
 (a;q)_\infty=\prod_{k=0}^{\infty}(1-aq^k).
\end{equation}
For each fixed \(\mu>0\), the ratio
\(\kappa_q(\mu+n)/(1-q)^n\) is bounded above and below by
positive constants independent of \(n\). Consequently,
\eqref{ob:mass} is precisely the convergence condition for the
nonnegative series in~\eqref{ob:series}, apart from the requirement
that the sum be nonzero.

\subsection{An exact variance identity}
\label{ob:variance-subsection}

Let \(\psi_q(t)=(\log\Gamma_q)'(t)\). Differentiating the product
in~\eqref{ob:reciprocal-product} gives, for \(t>0\),
\begin{align}
 \psi_q(t)
 &=-\log(1-q)+(\log q)\sum_{k=0}^{\infty}
       \frac{q^{t+k}}{1-q^{t+k}},\label{ob:psi}\\
 \psi_q'(t)
 &=(\log q)^2\sum_{k=0}^{\infty}
       \frac{q^{t+k}}{(1-q^{t+k})^2}>0.
       \label{ob:psi-prime}
\end{align}
For \(\mu>0\), define a probability distribution on the
nonnegative integers by
\begin{equation}\label{ob:weights}
 p_n(\mu,z)=\frac{g_nz^n\kappa_q(\mu+n)}{G(\mu;z)}.
\end{equation}
Expectation and variance under this distribution will be denoted by
\(\mathbb E_{\mu,z}\) and \(\operatorname{Var}_{\mu,z}\), and
\(\mathsf N\) denotes its integer-valued random variable.

\begin{proposition}[Curvature as variance minus mean curvature]
\label{ob:variance}
For any nonnegative coefficient sequence satisfying
\eqref{ob:mass}, the function \(G(\mu;z)\) is positive and
twice continuously differentiable for \(\mu>0\), and
\begin{equation}\label{ob:identity}
 \frac{\mathrm d^2}{\mathrm d\mu^2}\log G(\mu;z)
 =\operatorname{Var}_{\mu,z}\!\bigl(\psi_q(\mu+\mathsf N)\bigr)
  -\mathbb E_{\mu,z}\!\bigl[\psi_q'(\mu+\mathsf N)\bigr].
\end{equation}
Log-concavity or the absence of internal zeros in \((g_n)\) is not
needed for this identity.
\end{proposition}

\begin{proof}
On every compact parameter interval \([\eta,L]\subset(0,\infty)\),
the functions \(\psi_q(\mu+n)\) and \(\psi_q'(\mu+n)\) are
uniformly bounded in \(\mu\) and \(n\). This follows from
\eqref{ob:psi}--\eqref{ob:psi-prime}, since their denominators are
bounded away from zero and their numerators admit summable geometric
bounds for \(\mu+n\ge\eta\). Also,
\[
 0<\kappa_q(\mu+n)
 \le\frac{(1-q)^{\eta-1}}{(q;q)_\infty}(1-q)^n.
\]
Because
\[
 \kappa_q'(t)=-\psi_q(t)\kappa_q(t),\qquad
 \kappa_q''(t)=\bigl(\psi_q(t)^2-\psi_q'(t)\bigr)\kappa_q(t),
\]
condition~\eqref{ob:mass} uniformly dominates the original series
and its first two parameter derivatives on the compact interval.
Termwise differentiation is therefore valid. Dividing the first two
derivatives by \(G\) gives
\[
 \frac{G'}G=-\mathbb E_{\mu,z}[\psi_q(\mu+\mathsf N)],
 \qquad
 \frac{G''}G=
 \mathbb E_{\mu,z}[\psi_q(\mu+\mathsf N)^2-\psi_q'(\mu+\mathsf N)].
\]
Subtracting \((G'/G)^2\) proves~\eqref{ob:identity}.
\end{proof}

\paragraph{The unproved inequality.}
The specific bridge still required is
\begin{equation}\label{ob:target}
 \boxed{\quad
 \operatorname{Var}_{\mu,z}\!\bigl(\psi_q(\mu+\mathsf N)\bigr)
 \le \mathbb E_{\mu,z}\!\bigl[\psi_q'(\mu+\mathsf N)\bigr]
 \quad}
\end{equation}
for every \(\mu>0\), every \(0<q<1\), and every doubly positive
\((g_n)\) and \(z\) satisfying~\eqref{ob:mass}. Equation
\eqref{ob:target} is a proposed additional inequality, not a
consequence established by the preceding determinant arguments.

\begin{proposition}[What the bridge would imply]
\label{ob:turan-equivalence}
For fixed admissible \(q,(g_n),z\), inequality~\eqref{ob:target}
for all \(\mu>0\) is equivalent to logarithmic concavity of
\(G(\cdot;z)\) on \((0,\infty)\). It is also equivalent to
\begin{equation}\label{ob:turan}
 \begin{split}
 T_{\alpha,\beta}(\mu;z)
 &:={G(\mu+\alpha;z)G(\mu+\beta;z)}\\
 &\hspace{1.2em}-{G(\mu;z)G(\mu+\alpha+\beta;z)}\ge0
 \end{split}
\end{equation}
for all real \(\alpha,\beta\ge0\) and \(\mu>0\).
If these equivalent assertions hold, \eqref{ob:turan} also holds
at \(\mu=0\), with the reciprocal \(q\)-gamma function extended
continuously by \(\kappa_q(0)=0\).
\end{proposition}

\begin{proof}
The first equivalence is~\eqref{ob:identity}. Put
\(L(\mu)=\log G(\mu;z)\). For \(\mu>0\),
\begin{align*}
 &L(\mu+\alpha)+L(\mu+\beta)-L(\mu)-L(\mu+\alpha+\beta)\\
 &\qquad=-\int_0^\alpha\int_0^\beta
 L''(\mu+s+t)\,\mathrm dt\,\mathrm ds.
\end{align*}
Thus \(L''\le0\) implies~\eqref{ob:turan}. Conversely, setting
\(\alpha=\beta=h>0\), taking logarithms, dividing by \(h^2\),
and letting \(h\downarrow0\) gives \(L''(\mu)\le0\).

For the endpoint, formula~\eqref{ob:reciprocal-product} extends
\(\kappa_q\) continuously to zero, and for every \(\mu\ge0\)
\[
 0\le\kappa_q(\mu+n)
 \le\frac{(1-q)^{n-1}}{(q;q)_\infty}.
\]
Dominated convergence using~\eqref{ob:mass} makes \(G\)
continuous at \(\mu=0\). Taking the limit in
\eqref{ob:turan} proves the endpoint assertion, even when
\(G(0;z)=0\); no logarithm at that endpoint is required.
\end{proof}

The target is an inequality for the value of the Tur\'anian at
admissible \(z\). It does not assert nonnegativity of each power
series coefficient, nor does it assert a result at negative \(\mu\).

\subsection{What the published conjecture and corollary establish}
\label{ob:published-status}

For the series considered here, \cite[Theorem~3]{KK2018} proves a
function-level Tur\'an inequality for positive integer \(\alpha\),
real \(\beta>0\), and \(\mu\ge\max(-\beta,-1)\); its coefficientwise
conclusion has an additional restriction. Conjecture~1 removes the
shift restriction and includes negative \(\mu\). Those broader
claims are not established here.

Corollary~3.1 of that paper claims multiplicative concavity of
\begin{equation}\label{ob:ghat}
 \widehat G_q(a;y)=(a;q)_\infty
 \sum_{n=0}^{\infty}\frac{g_ny^n}{(a;q)_n},\qquad 0<a<1.
\end{equation}
The identity
\begin{equation}\label{ob:ghat-relation}
 G(\mu;z)=\frac{(1-q)^{\mu-1}}{(q;q)_\infty}
 \widehat G_q\bigl(q^\mu;(1-q)z\bigr)
\end{equation}
shows why that conclusion would imply the positive-parameter target.
Its displayed proof obtains
\begin{equation}\label{ob:fixed-step}
 \widehat G_q(a;y)^2\ge
 \widehat G_q(a/q;y)\widehat G_q(aq;y)
\end{equation}
for admissible arguments and then invokes its Lemma~1 (p.~334).
That lemma requires every dilation of one fixed function, whereas
\eqref{ob:fixed-step} supplies the fixed dilation \(q\).
Changing the base also changes \(\widehat G_q\), so it cannot
supply the missing quantifier. This is a gap in the cited derivation,
not a disproof of its conclusion. No exhaustive later-literature or
publication-priority claim is made. The unproved step in the present
paper remains \eqref{ob:target}.

\begin{example}[All-order signs do not supply nonlinear closure]
\label{ob:gaussian}
The limitation is already visible for an elementary kernel. On
positive sums, let
\[
 \widetilde K(x,y)=e^{-4(x+y)^2}
 =e^{-4x^2}e^{-4y^2}e^{-8xy}.
\]
Positive row and column factors and the exponential determinant show
that this kernel is \(\SSR_\infty\) with signature
\((-1)^{\binom r2}\). The finite sequence
\((1,e^{12},0,0,\ldots)\) is doubly positive. Nevertheless,
\[
 H(\mu)=e^{-4\mu^2}+e^{12}e^{-4(\mu+1)^2}
\]
is not log-concave on \((0,\infty)\). At \(\mu=1\) its two
summands are equal; their logarithmic derivatives are \(-8\) and
\(-16\), and their logarithmic second derivatives both equal
\(-8\). Hence
\[
 (\log H)''(1)=-8+\frac{(-8+16)^2}{4}=8>0.
\]
This is a counterexample to an abstract inference from sign regularity
and doubly positive coefficients. It is not a counterexample to the
reciprocal \(q\)-gamma inequality~\eqref{ob:target}.
\end{example}

\section{An arithmetic limitation of determinant signs}
\label{ob:arithmetic}

The all-order signs also do not determine the arithmetic nature of
special values of derivatives. The following explicit invariance
makes the limitation precise.

\begin{proposition}[Exponential normalization preserves every sign]
\label{ob:gauge}
For any \(a\in\R\), set
\[
 f_a(s)=\frac{e^{a(s-1)}}{\Gamma(s)},\qquad s>0.
\]
The kernel \(f_a(x+y)\) has exactly the same ordered determinant
signs as \(1/\Gamma(x+y)\). Moreover,
\begin{equation}\label{ob:gauge-derivative}
 f_a(1)=1,\qquad f_a'(1)=\gamma+a,
\end{equation}
where \(\gamma\) is Euler's constant.
\end{proposition}

\begin{proof}
For increasing row and column nodes with positive sums,
\[
 \det[f_a(x_i+y_j)]_{i,j=1}^r
 =\exp\!\left(a\left(\sum_{i=1}^r x_i+
                       \sum_{j=1}^r y_j-r\right)\right)
   \det\!\left[\frac1{\Gamma(x_i+y_j)}\right]_{i,j=1}^r.
\]
The prefactor is strictly positive. The normalization follows from
\(\Gamma(1)=1\); differentiation and
\(\psi(1)=-\gamma\) give~\eqref{ob:gauge-derivative}, where
\(\psi=\Gamma'/\Gamma\); see \cite[Section~5.4(ii)]{DLMF}.
\end{proof}

As \(a\) ranges over \(\R\), the derivative in
\eqref{ob:gauge-derivative} ranges over all real numbers while the
normalization and every determinant sign remain unchanged. Those
data alone therefore cannot force that derivative to be rational,
irrational, algebraic, or transcendental. No arithmetic classification
of Euler's constant follows from the sign-regularity results proved
here; such a conclusion would require additional arithmetic input.


\section{Final scope}

The determinant classifications are global on the stated
positive-argument domains. Discrete and integral composition supply
strictness, while the applications retain their rank and convergence
conditions. The variance bound in Section~\ref{ob:nonlinear} remains
an additional nonlinear target; none of the proved determinant or
shape results depends on it.

\begin{thebibliography}{9}
\bibitem{KVZ}
D. Karp, A. Vishnyakova, and Y. Zhang,
\emph{Unimodality preservation by ratios of functional series and
integral transforms}, Results in Mathematics \textbf{80}, 112 (2025).
\href{https://doi.org/10.1007/s00025-025-02431-4}{doi:10.1007/s00025-025-02431-4}.
The comparisons use Theorem~2.5 and Section~3, Examples~6, 8, and~11,
of \href{https://arxiv.org/abs/2408.01755v2}{arXiv:2408.01755v2}
(1 March 2025).

\bibitem{KK2018}
S. I. Kalmykov and D. B. Karp,
\emph{Inequalities for series in \(q\)-shifted factorials and
\(q\)-gamma functions}, Journal of Mathematical Analysis and
Applications \textbf{460} (2018), 332--351.
\href{https://doi.org/10.1016/j.jmaa.2017.11.041}{doi:10.1016/j.jmaa.2017.11.041}.
The source comparison concerns Lemma~1 (p.~334), Theorem~3,
Conjecture~1, and Corollary~3.1 (p.~344).

\bibitem{DLMF}
NIST Digital Library of Mathematical Functions,
\href{https://dlmf.nist.gov/5.2.E1}{Euler's gamma integral, Eq.~5.2.1};
\href{https://dlmf.nist.gov/5.12.E1}{Euler's beta integral, Eq.~5.12.1};
\href{https://dlmf.nist.gov/5.18}{\(q\)-gamma functions, Section~5.18};
and \href{https://dlmf.nist.gov/5.4.ii}{special values of the
psi function, Section~5.4(ii)}. Accessed 20 November 2026.
\end{thebibliography}

\end{document}
